Another one of our models that we experimented with remains a \textbf{logistic regression} classifier, but extends the TF--IDF baseline by combining three feature types extracted from each lyric. First, we compute \textbf{word-level TF--IDF} features using unigrams and bigrams, removing English stopwords and pruning extremely rare/common terms (\texttt{min\_df=2}, \texttt{max\_df=0.9}), with sublinear TF scaling to reduce the effect of repetition. Second, we add \textbf{character-level TF--IDF} features using 3--6 character n-grams (\texttt{char\_wb}) to capture stylistic patterns such as orthographic variation and elongated spellings that may be missed by word tokenization. Third, we include a small set of \textbf{numeric lyric-structure features} (e.g., word/line counts, bracket-tag count, punctuation, capitalization, and type--token ratio), standardized and concatenated with the sparse TF--IDF vectors. The resulting combined feature vector is used to train logistic regression with L2 regularization (optimized with \texttt{lbfgs}); we tuned the inverse regularization strength and selected $C=8$.

On the balanced validation set, the \textbf{majority-class baseline} achieved \textbf{Accuracy = 0.50} and \textbf{Macro-F1 = 0.33}, while the \textbf{trained TF--IDF + logistic regression baseline} achieved \textbf{Accuracy = 0.85} and \textbf{Macro-F1 = 0.85}. Our proposed feature-engineered model achieved \textbf{Accuracy = 0.873} and \textbf{Macro-F1 = 0.873}, improving over the majority baseline by \textbf{+0.373} accuracy and \textbf{+0.543} macro-F1, and improving over the trained TF--IDF baseline by approximately \textbf{+0.023} on both metrics.